% =====================================================================
%  config/commands.tex
%  ---------------------------------------------------------------
%  Sistema de marcado de la plantilla. Define CINCO tipos de bloque
%  visualmente distinguibles, de modo que quien lea el documento
%  identifique de inmediato la naturaleza de cada texto:
%
%    instruccion        -> qué debe escribir el equipo (guía de la plantilla)
%    ejemplo            -> ilustración; NO es contenido a conservar
%    politicaacademica  -> regla local del curso, NO requisito de la norma
%    notanormativa      -> referencia o aclaración conceptual con fuente
%    observacion        -> origen de una corrección respecto de la plantilla
%                          anterior; SÓLO aparece en la variante anotada
%
%  Requiere que config/metadata.tex se cargue ANTES que este archivo,
%  porque aquí se consultan los interruptores \ifmostrarinstrucciones
%  y \ifmostrarejemplos.
% =====================================================================

% ---------------------------------------------------------------------
% Estilo base común a todas las cajas
% ---------------------------------------------------------------------
\tcbset{
  cajaplantilla/.style={
    enhanced, breakable,
    boxrule=0.6pt, arc=2pt,
    left=6pt, right=6pt, top=5pt, bottom=5pt,
    fonttitle=\bfseries\footnotesize,
    fontupper=\small,
    attach boxed title to top left={xshift=6pt, yshift=-2pt},
    boxed title style={boxrule=0pt, arc=1pt, left=4pt, right=4pt, top=1pt, bottom=1pt},
    toptitle=2pt,
    before skip=8pt, after skip=8pt
  }
}

% ---------------------------------------------------------------------
% 1. INSTRUCCIÓN — guía de la plantilla (se oculta en la versión final)
% ---------------------------------------------------------------------
\ifmostrarinstrucciones
  \newtcolorbox{instruccion}[1][Qué debe documentarse en este apartado]{%
    cajaplantilla,
    colback=instrfondo, colframe=instrverde,
    coltitle=white, colbacktitle=instrverde,
    title={#1}}
\else
  \excludecomment{instruccion}
\fi

% ---------------------------------------------------------------------
% 2. EJEMPLO — ilustrativo, se elimina en la versión final
% ---------------------------------------------------------------------
\ifmostrarejemplos
  \newtcolorbox{ejemplo}[1][Ejemplo ilustrativo]{%
    cajaplantilla,
    colback=ejemfondo, colframe=ejemazul,
    coltitle=white, colbacktitle=ejemazul,
    title={#1}}
\else
  \excludecomment{ejemplo}
\fi

% ---------------------------------------------------------------------
% 3. POLÍTICA ACADÉMICA — regla local del curso. SIEMPRE visible.
%    Se conserva porque el curso la exige, pero queda marcada para que
%    nunca se confunda con un requisito de ISO/IEC/IEEE 29119.
% ---------------------------------------------------------------------
\newtcolorbox{politicaacademica}[1][Política académica de la Experiencia Educativa]{%
  cajaplantilla,
  colback=polfondo, colframe=polinaranja,
  coltitle=white, colbacktitle=polinaranja,
  title={#1}}

% ---------------------------------------------------------------------
% 4. NOTA NORMATIVA / CONCEPTUAL — con fuente identificable
% ---------------------------------------------------------------------
\newtcolorbox{notanormativa}[1][Referencia normativa]{%
  cajaplantilla,
  colback=normfondo, colframe=normgris,
  coltitle=white, colbacktitle=normgris,
  title={#1}}

% ---------------------------------------------------------------------
% 5. OBSERVACIÓN DE LA REVISIÓN — trazabilidad del cambio.
%    Sólo se compone en la variante anotada (make annotated). No forma
%    parte del plan: documenta por qué la plantilla es como es.
%
%    \begin{observacion}[Título corto]
%      \obsproblema{Qué estaba mal en la plantilla anterior.}
%      \obscambio{Qué se hizo en su lugar.}
%      \obsjustificacion{Por qué ese cambio resuelve el problema.}
%      \obsfuente{OBS §4.4}
%    \end{observacion}
%
%    Las claves de \obsfuente siguen el criterio de docs/MATRIZ_CAMBIOS.md:
%      OBS §n  -> Observaciones_plantilla_pruebas_ISO_29119.docx, apartado n
%      PDF §n  -> Plan_de_trabajo_29119_AngelAguilar.pdf, sección n
%      DIAG    -> análisis de los diagramas de referencia
%      PROPIA  -> criterio de esta reingeniería
% ---------------------------------------------------------------------
\newcommand{\obsitem}[2]{%
  \par\addvspace{2pt}%
  \noindent\textbf{\textcolor{obsmorado}{#1}}\enspace #2\par}
\newcommand{\obsproblema}[1]{\obsitem{Problema identificado.}{#1}}
\newcommand{\obscambio}[1]{\obsitem{Cambio aplicado.}{#1}}
\newcommand{\obsjustificacion}[1]{\obsitem{Justificación.}{#1}}
\newcommand{\obsfuente}[1]{\obsitem{Fuente.}{#1\ \textperiodcentered\ detalle completo en \texttt{docs/MATRIZ\_CAMBIOS.md}}}

\ifmostrarobservaciones
  \newtcolorbox{observacion}[1][]{%
    cajaplantilla,
    colback=obsfondo, colframe=obsmorado,
    coltitle=white, colbacktitle=obsmorado,
    title={Observación de la revisión\ifstrempty{#1}{}{: #1}}}
\else
  \excludecomment{observacion}
\fi

% ---------------------------------------------------------------------
% Menciones a las figuras
%   Los diagramas son material de apoyo y no se componen en la variante
%   final. \sidiag envuelve las frases (o los incisos) que remiten a una
%   figura, de modo que el texto siga siendo correcto cuando no está.
%   Redacte siempre la frase de manera que se sostenga sin el inciso.
% ---------------------------------------------------------------------
\newcommand{\sidiag}[1]{\ifmostrardiagramas #1\fi}

% ---------------------------------------------------------------------
% Marcador de campo por rellenar
% ---------------------------------------------------------------------
% Se compone como texto normal (no como caja) para que pueda partirse
% entre líneas dentro de las celdas estrechas de las tablas. El color y
% los corchetes bastan para localizarlo; al sustituirlo por contenido
% real, el marcador desaparece por completo.
\newcommand{\campo}[1]{{\color{campotexto}\itshape[\hspace{0pt}#1\hspace{0pt}]}}

% Marcador breve dentro de tablas (sin caja, para no romper filas)
\newcommand{\pordef}[1][\ ]{\textit{\textcolor{gray}{#1}}}

% ---------------------------------------------------------------------
% Divisor de PARTE (PARTE I / PARTE II)
% ---------------------------------------------------------------------
\newcommand{\parteplan}[2]{%
  \clearpage
  \phantomsection
  \addcontentsline{toc}{part}{#1. #2}%
  \vspace*{5\baselineskip}
  \begin{center}
    {\Large\bfseries\color{uvazul} #1}\\[10pt]
    {\LARGE\bfseries\color{uvazul} #2}
  \end{center}
  \vspace{2\baselineskip}
  \setcounter{section}{0}%
  \clearpage
}

% ---------------------------------------------------------------------
% Inicio de los anexos: la numeración pasa de 1,2,3 a ANEXO A, B, C…
% ---------------------------------------------------------------------
\newcommand{\iniciaranexos}{%
  \clearpage
  \setcounter{section}{0}%
  \renewcommand{\thesection}{ANEXO \Alph{section}}%
  \renewcommand{\thesubsection}{\Alph{section}.\arabic{subsection}}%
  % ensancha la caja del número en el índice sólo a partir de los anexos
  \addtocontents{toc}{\protect\cftsetindents{section}{1.5em}{6.4em}}%
  \phantomsection
  \addcontentsline{toc}{part}{ANEXOS}%
}

% ---------------------------------------------------------------------
% Nota al pie de una tabla (fuente / aclaración)
% ---------------------------------------------------------------------
\newcommand{\notatabla}[1]{\par\vspace{3pt}{\footnotesize\itshape #1\par}}

% ---------------------------------------------------------------------
% Encabezado de tabla sombreado
% ---------------------------------------------------------------------
\newcommand{\filaenc}{\rowcolor{uvgris}}

% ---------------------------------------------------------------------
% Abreviatura de la norma, para citarla siempre igual
% ---------------------------------------------------------------------
\newcommand{\norma}{ISO/IEC/IEEE 29119}
\newcommand{\normap}[1]{ISO/IEC/IEEE 29119-#1}
